%==============================================================================%
% Chapter 8 – Algebraic NTT Infrastructure (NTTFullAlgebra.ec)
%==============================================================================%

\chapter{Algebraic NTT Infrastructure}
\label{ch:algebra}

\section{Overview of \texttt{NTTFullSpec.ec} and \texttt{NTTFullAlgebra.ec}}

The abstract spectral NTT specification lives in two files:
\begin{itemize}
  \item \texttt{NTTFullSpec.ec} (160 lines) defines the abstract evaluation
        maps \ec{full\_ntt} and \ec{full\_invntt} and states the main
        correctness lemmas.
  \item \texttt{NTTFullAlgebra.ec} (1000+ lines) develops the algebraic
        machinery needed to prove those lemmas: primitive root properties,
        bit-reversal identities, and butterfly-stage decompositions.
\end{itemize}

\section{Primitive Root Properties}

\subsection{Existence and Uniqueness}

The prime $q = 64513$ satisfies $q \equiv 1 \pmod{512}$.  Therefore $\F_q^*$
contains an element of order $512$, i.e., a primitive $512$th root of unity.

\begin{theorem}[Primitive root existence]
  The element $\zeta_0 = 426 \in \F_q$ has multiplicative order exactly $512$.
\label{thm:zeta_order}
\end{theorem}

\begin{proof}
  We verify:
  \begin{enumerate}
    \item $\zeta_0^{512} \equiv 1 \pmod{64513}$ (direct computation).
    \item For all $d \mid 512$ with $d < 512$:
          $\zeta_0^d \not\equiv 1 \pmod{64513}$.
          It suffices to check $d = 256$ and $d = 128$ (since $512 = 2^9$
          and the order must divide $512$).
  \end{enumerate}
  Both computations are discharged by the \ec{smt} tactic.
\end{proof}

\begin{corollary}
  $\zeta_0^{256} \equiv -1 \pmod{q}$ and $\zeta_0^{128} \not\equiv \pm 1 \pmod{q}$.
\label{cor:zeta256}
\end{corollary}

\subsection{Orthogonality}

The $512$th roots of unity in $\F_q$ are $\{\zeta_0^k : k = 0, 1, \dots, 511\}$.
The primitive roots are those $\zeta_0^k$ with $\gcd(k, 512) = 1$.  The NTT
formula uses the set of \emph{odd} powers $\{\zeta_0^{2i+1} : i = 0, \dots, 255\}$,
which are exactly the primitive roots.

\begin{lemma}[Orthogonality relation]
  For $i, j \in \{0, \dots, 255\}$ with $i \ne j$:
  \[
    \sum_{k=0}^{255} \zeta_0^{(2i+1)k} \cdot \zeta_0^{-(2j+1)k} = 0 \in \F_q.
  \]
  For $i = j$, the sum equals $256$.
\label{lem:orthogonality}
\end{lemma}

\begin{proof}
  For $i \ne j$: the sum is a geometric series with ratio
  $r = \zeta_0^{2(i-j)}$.  Since $\zeta_0$ has order 512, $r$ is a
  $256$th root of unity.  The exponent $2(i-j) \not\equiv 0 \pmod{256}$
  (because $|i-j| < 256$, so $2|i-j| < 512$, and $2(i-j) \ne 0$).
  Thus $r^{256} = \zeta_0^{512(i-j)} = 1$ and $r \ne 1$, so
  $\sum_{k=0}^{255} r^k = (r^{256} - 1)/(r-1) = 0$.

  For $i = j$: each term is $\zeta_0^0 = 1$, and there are 256 terms.
\end{proof}

\section{Bit-Reversal Permutation}
\label{sec:brev}

\subsection{Definition}

The bit-reversal permutation $\brev_8 : \{0,\dots,255\} \to \{0,\dots,255\}$
reverses the 8 binary digits of its input:
\[
  \brev_8(b_7 b_6 b_5 b_4 b_3 b_2 b_1 b_0) = b_0 b_1 b_2 b_3 b_4 b_5 b_6 b_7.
\]

\begin{lstlisting}[language=EasyCrypt,
  caption={Bit-reversal definition in \texttt{NTTFullAlgebra.ec}},
  label={lst:brev}]
op bsrev (n k : int) : int =
  bsrev8 k.   (* 8-bit bit reversal *)

op bsrev8 (k : int) : int =
  let k0 = k %% 2 in
  let k1 = (k %/ 2) %% 2 in
  (* ... *)
  let k7 = (k %/ 128) %% 2 in
  k0 * 128 + k1 * 64 + k2 * 32 + k3 * 16 +
  k4 * 8   + k5 * 4  + k6 * 2  + k7.
\end{lstlisting}

\subsection{Key Properties}

\begin{lemma}[Bit-reversal involutivity]
  $\brev_8(\brev_8(k)) = k$ for all $k \in [0, 256)$.
\label{lem:brev_invol}
\end{lemma}

\begin{lemma}[Bit-reversal and halving]
  For $k \in [1, 256)$ and $\ell = 2^s$ (a power of 2 with $s \le 7$):
  \[
    \brev_8(k \bmod \ell) = \brev_8(k) / (256 / \ell) \cdot (256 / \ell) + \brev_8(k) \bmod (256 / \ell).
  \]
\label{lem:brev_halve}
\end{lemma}

This lemma, though technical, is essential for connecting the loop index $k$
to the spectral index $\brev_8(k)$.

\subsection{Bit-Reversal and Twiddle Factor Indexing}

\begin{lemma}[Twiddle factor ordering]
  In the forward NTT outer loop with half-length $\ell = 2^{7-s}$, after
  completing $g$ groups of stage $s$, the twiddle counter satisfies
  $k = 2^s + g$, and the twiddle factor is $\zeta_0^{\brev_8(2^s + g)}$.
\label{lem:twiddle_index}
\end{lemma}

\begin{proof}
  By induction on $s$ (outer loops) and $g$ (groups within a stage).
  The key observation is that as the outer counter increments,
  the bit-reversal of the counter gives exactly the correct spectral index
  for the butterfly group.
\end{proof}

\section{Stage Decomposition Theorem}

\subsection{Partial NTT}

\begin{definition}[Stage-$s$ partial NTT]
  For $0 \le s \le 8$, the stage-$s$ partial NTT of a vector
  $f \in \F_q^{256}$ is defined by:
  \[
    \NTT_s(f)[i] = \sum_{j \equiv i \pmod{2^s}} f[j] \cdot
    \zeta_0^{\brev_8(\floor{i/2^s}) \cdot j},\quad i = 0, \dots, 255.
  \]
\label{def:partial_ntt}
\end{definition}

The full NTT corresponds to $\NTT_8$ (all 8 stages completed).  The identity
$\NTT_0(f) = f$ holds trivially (the empty product is 1).

\subsection{Butterfly Stage Lemma}

\begin{lemma}[Single-stage butterfly correctness]
  For all $s \in \{0, \dots, 7\}$ and $f \in \F_q^{256}$:
  \[
    \NTT_{s+1}(f) = \op{butterfly\_stage}(s, \NTT_s(f)),
  \]
  where $\op{butterfly\_stage}(s, a)$ applies one complete pass of butterflies
  with half-length $\ell = 2^{7-s}$.
\label{lem:butterfly_stage}
\end{lemma}

\begin{proof}
  Fix an index $i$ and let $g = \floor{i / 2^{s+1}}$ (the group number),
  $r = i \bmod 2^{s+1}$ (the position within the group),
  $r_0 = r \bmod 2^s$ (the low-half position).

  The butterfly for group $g$ uses twiddle factor $w = \zeta_0^{\brev_8(2^s + g)}$.

  Case $r < 2^s$ (upper butterfly input):
  \begin{align*}
    \op{butterfly\_stage}(s, a)[i]
    &= a[i] + w \cdot a[i + 2^s] \\
    &= \sum_{j \equiv r_0 \pmod{2^s}} f[j] \zeta_0^{\brev_8(g) j}
     + \zeta_0^{\brev_8(2^s+g)} \sum_{j \equiv r_0 \pmod{2^s}} f[j+2^s] \zeta_0^{\brev_8(g) j}.
  \end{align*}
  (By the inductive hypothesis $\NTT_s(f)[i] = \sum_{j \equiv i \pmod{2^s}} f[j]
  \zeta_0^{\brev_8(\floor{i/2^s})j}$.)

  Combining the sums over $j$ and $j + 2^s$, and using the identity
  $\brev_8(2^s + g) = \brev_8(g)/2 + 2^{7-s} \cdot (g \bmod 2^s)$
  (from \Cref{lem:brev_halve}), one can verify that the result equals
  $\NTT_{s+1}(f)[i]$ by direct algebraic manipulation.

  The case $r \ge 2^s$ is symmetric.  The full algebraic proof in
  \easycrypt{} uses the \ec{ring} tactic after reducing to finite sums.
\end{proof}

\subsection{Full NTT Correctness}

\begin{theorem}[Forward NTT correctness]
  The imperative procedure \ec{NTT.ntt} computes the full NTT:
  \[
    \forall a \in \Z_q^{256}.\;
    \ec{NTT.ntt}(a) = \NTT_8(a) = \op{full\_ntt}(a).
  \]
\label{thm:ntt_correct}
\end{theorem}

\begin{proof}
  By induction on the outer loop iteration, using \Cref{lem:butterfly_stage}
  as the inductive step.  The loop invariant asserts that after $s$ outer
  loop iterations, the array equals $\NTT_s$ of the input.

  The base case ($s = 0$): the array is unchanged, matching $\NTT_0 = \op{id}$.
  The inductive step: applying the butterfly stage transforms $\NTT_s$ to
  $\NTT_{s+1}$ by \Cref{lem:butterfly_stage}.
  The final case ($s = 8$): $\NTT_8 = \op{full\_ntt}$ by \Cref{def:full_ntt}.
\end{proof}

\section{Inverse NTT Algebra}

\subsection{Abstract Inverse NTT}

\begin{definition}[Abstract inverse NTT, \ec{full\_invntt}]
  \[
    \INTT(p)[i] = 256^{-1} \sum_{j=0}^{255} p[j] \cdot \zeta_0^{-((2\brev_8(j)+1) \cdot i)}
  \]
\label{def:full_intt2}
\end{definition}

\begin{theorem}[Inverse NTT correctness]
  $\op{full\_invntt}(\op{full\_ntt}(p)) = p$ for all $p \in \Z_q^{256}$.
\label{thm:intt_correct}
\end{theorem}

\begin{proof}
  Substituting the definition of $\op{full\_ntt}$:
  \begin{align*}
    \INTT(\NTT(p))[i]
    &= 256^{-1} \sum_{j=0}^{255} \NTT(p)[j] \cdot \zeta_0^{-((2\brev_8(j)+1) i)} \\
    &= 256^{-1} \sum_{j=0}^{255}
       \left(\sum_{k=0}^{255} p[k] \cdot \zeta_0^{(2\brev_8(j)+1)k}\right)
       \cdot \zeta_0^{-((2\brev_8(j)+1)i)} \\
    &= 256^{-1} \sum_{k=0}^{255} p[k]
       \sum_{j=0}^{255} \zeta_0^{(2\brev_8(j)+1)(k-i)}.
  \end{align*}

  For $k = i$: the inner sum is $\sum_{j=0}^{255} 1 = 256$.
  For $k \ne i$: the inner sum is 0 by a variant of \Cref{lem:orthogonality}
  (substituting $m = \brev_8(j)$ and using that $\brev_8$ is a bijection on $[0,256)$).

  Thus $\INTT(\NTT(p))[i] = 256^{-1} \cdot p[i] \cdot 256 = p[i]$.
\end{proof}

\section{The Critical Difference from MLKEM}

\begin{remark}[HAETAE vs.\ MLKEM exponent structure]
  In MLKEM, the twiddle factor for position $(s, g)$ in the loop is
  $\zeta^{\brev_8(2g+1)}$ (with a $2g+1$ structure), which corresponds to
  the \emph{incomplete} $128$-point transform used in that scheme.  The
  algebra lemmas for MLKEM are stated in terms of $\brev_7$ (7-bit reversal)
  rather than $\brev_8$.

  In \haetae{}, the twiddle factor is $\zeta^{\brev_8(2^s + g)}$, where the
  bit-reversal acts on an index that includes the stage counter.  This produces
  the \emph{complete} $256$-point negacyclic NTT.  The corresponding algebra
  lemmas require $\brev_8$ (8-bit reversal) and the orthogonality relation over
  all $256$ primitive $512$th roots of unity.

  This distinction is the primary obstacle to reusing MLKEM's EasyCrypt
  proofs, and is the motivation for developing \texttt{NTTFullAlgebra.ec}
  from scratch.
\label{rem:haetae_mlkem}
\end{remark}
